\section*{Chapter \ref{ch:b2:wittig} --- Wittig and Other C=C Forming Reactions}

\begin{solution}{exo:b2:wittig:1}
\ce{Ph3P + CH3CH2Br -> [Ph3PCH2CH3]+ + Br-} (bimolecular substitution, heating in a solvent), then
butyllithium in dry ether or THF under inert gas: \ce{[Ph3PCH2CH3]+ + BuLi -> Ph3P=CHCH3 + BuH + Li+}.
\end{solution}

\begin{solution}{exo:b2:wittig:2}
The ylide \ce{Ph3P=CH2} and benzaldehyde give phenylethene (styrene), \ce{C6H5CH=CH2}, and
triphenylphosphine oxide.
\end{solution}

\begin{solution}{exo:b2:wittig:3}
\ce{Ph3P=CHCOOEt} and \ce{Ph3P=CHCOCH3}: stabilised. \ce{Ph3P=CHCH3}: non-stabilised, the one that
needs butyllithium (or sodium hydride, sodium amide). \ce{Ph3P=CHC6H5}: aryl-stabilised, made with an
alkoxide or hydroxide in practice, and giving mixtures of geometries.
\end{solution}

\begin{solution}{exo:b2:wittig:4}
Ethyl (\textit{E})-pent-2-enoate, \ce{CH3CH2CH=CHCOOEt}, with sodium diethyl phosphate.
\end{solution}

\begin{solution}{exo:b2:wittig:5}
Wittig: the ylide \ce{Ph3P=CH2}, made from methyltriphenylphosphonium iodide and butyllithium, with cyclohexanone gives only
methylenecyclohexane, the double bond outside the ring. Dehydration of 1-methylcyclohexan-1-ol goes
through a tertiary carbocation and loses the hydrogen that gives the more substituted alkene (Zaitsev):
mainly 1-methylcyclohexene, the wrong isomer.
\end{solution}

\begin{solution}{exo:b2:wittig:6}
Both cuts are the same: benzaldehyde and the benzylidene ylide \ce{Ph3P=CHC6H5}. That ylide is only
aryl-stabilised and gives an (\textit{E})/(\textit{Z}) mixture. For pure (\textit{E}), use the HWE
reaction of diethyl benzylphosphonate with benzaldehyde.
\end{solution}

\begin{solution}{exo:b2:wittig:7}
With \ce{Ph3P=CHCH3}: (\textit{Z})-pent-2-ene, \ce{CH3CH2CH=CHCH3}, mainly (\textit{Z}). With
\ce{Ph3P=CHCOOEt}: ethyl (\textit{E})-pent-2-enoate, mainly (\textit{E}).
\end{solution}

\begin{solution}{exo:b2:wittig:8}
\ce{Ph3P=CH2 + C6H10O -> C7H12 + Ph3PO}. Masses: methylenecyclohexane \qty{96.173}{g/mol},
triphenylphosphine oxide \qty{278.291}{g/mol}; atom economy $96.173/(96.173 + 278.291) =
96.173/374.464 \approx 25.7~\%$. Three quarters of the mass of the reagents leaves as phosphine oxide.
% ledger: aw:C, aw:H, aw:O, aw:P
\end{solution}

\begin{solution}{exo:b2:wittig:9}
Phosphorus, nucleophilic, displaces bromide:
\begin{equation*}
\ce{P(OEt)3 + BrCH2COOEt -> [(EtO)3PCH2COOEt]+ + Br-}.
\end{equation*}
Bromide
then attacks a carbon of an ethyl group of the phosphonium ion (bimolecular substitution), releasing
bromoethane and forming the \ce{P=O} bond:
\begin{equation*}
\ce{[(EtO)3PCH2COOEt]+ + Br- -> (EtO)2P(O)CH2COOEt + EtBr}.
\end{equation*}
\end{solution}

\begin{solution}{exo:b2:wittig:10}
\ce{OHC-CH=CH-CHO + 2Ph3P=CHC6H5 -> C6H5CH=CH-CH=CH-CH=CHC6H5 + 2Ph3PO}: two equivalents of the
benzylidene ylide (from benzyltriphenylphosphonium chloride and a base), one for each aldehyde group.
The new double bonds are formed as mixtures of geometries; the all-(\textit{E}) isomer, the most
stable, is isolated by crystallisation, or the mixture is isomerised.
\end{solution}

\begin{solution}{exo:b2:wittig:11}
\omterm{def:b2:enolates-aldol:aldol}{Aldol}: benzaldehyde with propanone in excess and sodium hydroxide (\cref{ch:b2:enolates-aldol});
cheap reagents, water as by-product, (\textit{E}) product, but the condensation can run twice
(dibenzalacetone) if propanone is not in excess. Wittig: benzaldehyde with the \omterm{def:b2:wittig:ylide}{stabilised ylide}
\ce{Ph3P=CHCOCH3}; (\textit{E}), one product, no double condensation, but the ylide is costly and
triphenylphosphine oxide must be removed. The \omterm{def:b2:enolates-aldol:aldol}{aldol} is the better route here.
\end{solution}

\begin{solution}{exo:b2:wittig:12}
Wittig: butanal with \ce{Ph3P=CHCH2CH2CH3} (from 1-bromobutane, then butyllithium), salt-free:
mainly (\textit{Z}); by-product triphenylphosphine oxide, and some (\textit{E}) isomer. Alkyne:
oct-4-yne (from ethyne by two alkylations with 1-bromopropane and sodium amide) hydrogenated over a
poisoned palladium catalyst (\cref{prop:b2:alkene-redox:syn-hydrogenation}): (\textit{Z}) only; the
by-products are sodium bromide and ammonia from the alkylations.
\end{solution}

\begin{solution}{pb:b2:wittig:1}
\textbf{1.} \ce{C23H46}, that is \ce{C_nH_{2n}} with $n = 23$: one double bond (or one ring), here a
\ce{C=C}.
\textbf{2.} The \ce{C9=C10} double bond.
\textbf{3.} Nonanal \ce{CH3(CH2)7CHO} with the ylide from a 1-halotetradecane; or tetradecanal
\ce{CH3(CH2)12CHO} with the ylide from a 1-halononane.
\textbf{4.} 1-Bromotetradecane, \ce{CH3(CH2)13Br}.
\textbf{5.} Bimolecular substitution by the bulky triphenylphosphine is fast at an unhindered primary
carbon and does not compete with elimination.
\textbf{6.} \ce{Ph3P + C14H29Br -> [Ph3PC14H29]+ + Br-}, tetradecyltriphenylphosphonium bromide.
\textbf{7.} Butyllithium (or sodium hydride, sodium amide); hydroxide is too weak a base for a
non-stabilised ylide, and the water it brings would protonate the ylide.
\textbf{8.} No, it carries only an alkyl group: mainly (\textit{Z}) (\cref{prop:b2:wittig:stereo}).
\textbf{9.} A four-membered ring \ce{P-C-C-O}: phosphorus bonded to the carbon carrying the
\ce{C13H27} chain, that carbon to the one carrying the \ce{C8H17} chain, and that one to oxygen, bonded
to phosphorus.
\textbf{10.} Butyllithium and the ylide react with water and oxygen.
\textbf{11.} The double bond joins the former carbonyl carbon of nonanal (C9 of the product) to the
former ylide carbon (C10), and no other bond moves (\cref{prop:b2:wittig:regiospecific}).
\textbf{12.} A carbocation at C9, which may lose a hydrogen from C8 or C10 and may shift: a mixture of
tricos-8-ene and tricos-9-ene, each as (\textit{E}) and (\textit{Z}), mostly (\textit{E}).
\textbf{13.} \ce{Ph3P=CHC13H27 + C8H17CHO -> C8H17CH=CHC13H27 + Ph3PO}.
\textbf{14.} Target \ce{C23H46}: \qty{322.621}{g/mol}; nonanal \ce{C9H18O}: \qty{142.242}{g/mol}; salt
\ce{C32H44BrP}: \qty{539.582}{g/mol}; \ce{Ph3PO}, \ce{C18H15OP}: \qty{278.291}{g/mol}.
\textbf{15.} The ylide \ce{C32H43P} weighs $539.582 - 80.912 = \qty{458.670}{g/mol}$; with nonanal,
\qty{600.912}{g/mol} of reagents for \qty{322.621}{g/mol} of product: $322.621/600.912 \approx
53.7~\%$.
\textbf{16.} By its polarity: the nonpolar hydrocarbon passes through a short silica column with a
nonpolar eluent while the oxide stays; or the oxide crystallises out of a cold nonpolar solvent.
\textbf{17.} The very strong \ce{P=O} bond formed outweighs the \ce{P-C} bond broken
(\cref{prop:b2:wittig:driving}).
\textbf{18.} It needs a \omterm{def:b2:wittig:hwe}{phosphonate} with an electron-withdrawing group on the reacting carbon, absent
from a plain alkyl chain, and it gives the (\textit{E})-alkene.
\textbf{19.} Tricos-9-yne, \ce{CH3(CH2)7C#C(CH2)12CH3}, with hydrogen over a poisoned palladium
catalyst (palladium on calcium carbonate with a lead salt and quinoline).
\textbf{20.} Dec-1-yne, \ce{CH3(CH2)7C#CH}, deprotonated by sodium amide, then 1-bromotridecane
\ce{CH3(CH2)12Br}: $8 + 2 + 13 = 23$ carbons, the triple bond between C9 and C10.
\textbf{21.} Further hydrogenation would give tricosane, with no double bond; the poisoned catalyst
binds the alkyne much more strongly than the alkene, which leaves the surface before a second addition.
\textbf{22.} Both take two steps from available pieces. The alkyne route gives the (\textit{Z}) isomer
cleanly and leaves only salts; the Wittig route gives some (\textit{E}) isomer and a mass of
triphenylphosphine oxide.
\textbf{23.} 100~\%: \ce{C23H44 + H2 -> C23H46}, every atom ends in the product.
\textbf{24.} $278.291/322.621 = \boldsymbol{\approx \qty{0.86}{g}}$ of triphenylphosphine oxide per
gram of pheromone.
\end{solution}
